\begin{definition}[The elementary binary tensor and its color differences]%
    \label{def:peps_kitaev_checkerboard_tensor}
    \lean{TNLean.PEPS.KitaevBit, TNLean.PEPS.KitaevBlockBoundary,
        TNLean.PEPS.KitaevBlockSpins, TNLean.PEPS.kitaevElementaryTensor,
        TNLean.PEPS.kitaevBlockColorSpins}
    \leanok
    All indices take values in $\mathbb Z_2$. Write the virtual legs in the
    order top, right, bottom, left. The two checkerboard orientations of the
    elementary tensor are
    \begin{align}
        T_0(t,r,d,l;k) & =\delta_{t,r}\delta_{d,l}\delta_{k,t+d}, &
        T_1(t,r,d,l;k) & =\delta_{t,l}\delta_{d,r}\delta_{k,t+d}.
    \end{align}
    These are the tensor of
    \cite[Section~7, ``Kitaev's code state'']{Schuch2010PEPS}
    and its rotated orientation. For four colors define
    \begin{align}
        F(p,q,r,s)=(p+q,q+r,r+s,s+p).
    \end{align}
    An open block has one pair of boundary labels on each side, ordered
    north, east, south, west. Each pair is read clockwise along its side.
    Its four physical spins are ordered clockwise from the upper-right site.
\end{definition}

\begin{definition}[The actual four-site checkerboard contraction]%
    \label{def:peps_kitaev_checkerboard_block}
    \lean{TNLean.PEPS.kitaevCheckerboardBlock, TNLean.PEPS.kitaevBoundaryCNOT}
    \leanok
    \uses{def:peps_kitaev_checkerboard_tensor}
    Let the boundary pairs be $N=(N_0,N_1)$, $E=(E_0,E_1)$,
    $S=(S_0,S_1)$ and $W=(W_0,W_1)$. Contract the four internal bonds of
    the $2\times2$ block around an A plaquette:
    \begin{align}
        B(N,E,S,W;\sigma)
        =\sum_{x_0,x_1,x_2,x_3}
         & T_0(N_0,x_0,x_3,W_1;\sigma_3)
        T_1(N_1,E_0,x_1,x_0;\sigma_0)\notag          \\
         & \qquad\cdot T_0(x_1,E_1,S_0,x_2;\sigma_1)
        T_1(x_3,x_2,S_1,W_0;\sigma_2).
    \end{align}
    % Source: SCP10 Section 7, paper_v3.tex:2758--2792; the explicit
    % coefficient formula above fixes all ports without a physical RG claim.
    % Ink: NW,NE,SE,SW carry T_0,T_1,T_0,T_1. Four virtual bonds x_0,...,x_3
    % are summed; N_0,N_1,E_0,E_1,S_0,S_1,W_0,W_1 stay open.
    % Physical NE,SE,SW,NW have sigma_0,sigma_1,sigma_2,sigma_3.
    % Boundary (virtual,physical)=(8,4). The internal A is not a tensor.
    \begin{tenkzequation}
        \tenkzkernel
        \begin{tenkz}[rows={wire,wire},cols=2,bonds=none]
            \tn[at=(1,1),name=blockTNW,skin=box,
                ports={90:virtual:$N_0$,0:virtual,270:virtual,180:virtual:$W_1$,
                        135:physical:$\sigma_3$}]{T_0}
            \tn[at=4 e of blockTNW,name=blockTNE,skin=box,
                ports={90:virtual:$N_1$,0:virtual:$E_0$,270:virtual,180:virtual,
                        45:physical:$\sigma_0$}]{T_1}
            \tn[at=3 s of blockTNE,name=blockTSE,skin=box,
                ports={90:virtual,0:virtual:$E_1$,270:virtual:$S_0$,180:virtual,
                        315:physical:$\sigma_1$}]{T_0}
            \tn[at=3 s of blockTNW,name=blockTSW,skin=box,
                ports={90:virtual,0:virtual,270:virtual:$S_1$,180:virtual:$W_0$,
                        225:physical:$\sigma_2$}]{T_1}
            \tnwire[name=blockX0]{blockTNW.0}{blockTNE.180}
            \tnwire[name=blockX1]{blockTNE.270}{blockTSE.90}
            \tnwire[name=blockX2]{blockTSE.180}{blockTSW.0}
            \tnwire[name=blockX3]{blockTSW.90}{blockTNW.270}
            \tnmark[form=label,label pos=n]{on blockX0 0.5}{$x_0$}
            \tnmark[form=label,label pos=e]{on blockX1 0.5}{$x_1$}
            \tnmark[form=label,label pos=s]{on blockX2 0.5}{$x_2$}
            \tnmark[form=label,label pos=w]{on blockX3 0.5}{$x_3$}
            \tnmark[form=label]{midway blockTNW and blockTSE}{A}
        \end{tenkz}
    \end{tenkzequation}
    The boundary transformation $C$ applies $(a,b)\mapsto(a,a+b)$ on
    each of the four pairs. This is the CNOT transformation in the
    blocking construction of \cite[Section~7]{Schuch2010PEPS}.
\end{definition}

\begin{lemma}[Boundary CNOTs]%
    \label{lem:peps_kitaev_checkerboard_cnot}
    \lean{TNLean.PEPS.kitaevBoundaryCNOT_apply,
        TNLean.PEPS.kitaevBoundaryCNOT_involutive}
    \leanok
    \uses{def:peps_kitaev_checkerboard_block}
    The boundary transformation acts on each pair by $(a,b)\mapsto(a,a+b)$
    and satisfies $C^2=1$.
\end{lemma}

\begin{proof}\leanok
    Addition is modulo two, so $a+(a+b)=b$.
\end{proof}

\begin{theorem}[Four-site blocking of Kitaev's elementary tensor]%
    \label{thm:peps_kitaev_checkerboard_blocking}
    \lean{TNLean.PEPS.kitaevCheckerboardBlock_apply,
        TNLean.PEPS.kitaevCheckerboardBlock_boundaryCNOT}
    \leanok
    \uses{def:peps_kitaev_checkerboard_block, lem:peps_kitaev_checkerboard_cnot}
    For every boundary labeling and every physical configuration,
    \begin{align}
        B(N,E,S,W;\sigma)
         & =\delta_{N_1,N_0}\delta_{E_1,E_0}
        \delta_{S_1,S_0}\delta_{W_1,W_0}
        \delta_{\sigma,F(N_0,E_0,S_0,W_0)}.
    \end{align}
    Consequently, if the four transformed boundary pairs are
    $(p,a_N),(q,a_E),(r,a_S),(s,a_W)$, then
    \begin{align}
        B\bigl(C((p,a_N),(q,a_E),(r,a_S),(s,a_W));\sigma\bigr)
         & =\delta_{a_N,0}\delta_{a_E,0}\delta_{a_S,0}\delta_{a_W,0}
        \delta_{\sigma,F(p,q,r,s)}.
    \end{align}
    Thus the actual contraction gives the color-difference tensor
    $K$ of \cite[Section~7, ``Kitaev's code state'']{Schuch2010PEPS}
    and four boundary registers fixed at zero. This is a coefficient
    identity for the open block; it does not assert a physical
    renormalization operation or a globally blocked lattice identity.
\end{theorem}

\begin{proof}\leanok
    A nonzero summand forces
    $x_0=N_0=N_1$, $x_1=E_0=E_1$, $x_2=S_0=S_1$ and $x_3=W_0=W_1$.
    There is therefore exactly one surviving internal labeling, and its
    four physical spins are $F(N_0,E_0,S_0,W_0)$. Substituting the boundary
    CNOTs gives the second identity, since $a+b=a$ exactly when $b=0$.
\end{proof}

\begin{definition}[Matrices of the blocked tensor and the retained colors]%
    \label{def:peps_kitaev_checkerboard_maps}
    \lean{TNLean.PEPS.kitaevCheckerboardBlockMatrix,
        TNLean.PEPS.kitaevBlockColorMatrix, TNLean.PEPS.kitaevBoundaryZeroEmbedding}
    \leanok
    \uses{def:peps_kitaev_checkerboard_block}
    Let $\mathsf B$ be the matrix of the actual four-site contraction, with
    physical rows and eight-bit boundary columns. Let $\mathsf K$ have
    entries $\mathsf K_{\sigma,p}=\delta_{\sigma,F(p)}$.
    Define the inclusion $E$ by fixing the four second boundary registers
    at zero:
    \begin{align}
        E\ket{p,q,r,s}=\ket{(p,0),(q,0),(r,0),(s,0)}.
    \end{align}
\end{definition}

\begin{theorem}[The actual tensor-map identity after the boundary unitary]%
    \label{thm:peps_kitaev_checkerboard_map_identity}
    \lean{TNLean.PEPS.kitaevBoundaryZeroEmbedding_apply,
        TNLean.PEPS.kitaevBoundaryCNOT_permMatrix_mem_unitaryGroup,
        TNLean.PEPS.kitaevBoundaryZeroEmbedding_isIsometry,
        TNLean.PEPS.kitaevCheckerboardBlockMatrix_mul_boundaryCNOT,
        TNLean.PEPS.kitaevBlockColorMatrix_eq_quantumDoubleDualTensor}
    \leanok
    \uses{def:peps_kitaev_checkerboard_maps,
        thm:peps_kitaev_checkerboard_blocking, def:peps_quantum_double_dual}
    The boundary CNOT matrix $\mathsf C$ is unitary, the inclusion satisfies
    $E^\dagger E=1$, and the actual blocked tensor map obeys
    \begin{align}
        \mathsf B\mathsf C=\mathsf K E^\dagger.
    \end{align}
    The color-difference tensor $\mathsf K$ is the dual toric-code tensor of
    Definition~\ref{def:peps_quantum_double_dual}: colors are read north,
    east, south, west. Its physical order is obtained by the cyclic reordering
    \[
        (\sigma_0,\sigma_1,\sigma_2,\sigma_3)
        \longmapsto(\sigma_3,\sigma_0,\sigma_1,\sigma_2).
    \]
    % Formula: B C = K E^dagger, with E fixing four virtual zero registers.
    % The K physical port denotes the unchanged tuple of four original spins.
    % Each zero* atom is a virtual covector <0|, with no physical output.
    % No bonds are contracted in this factorized expression; boundary=(8,1),
    % where the single physical port is the full four-spin tuple sigma.
    \begin{tenkzequation}
        $\mathsf B\mathsf C =$
        \tenkzkernel
        \begin{tenkz}[rows={wire,wire,wire},cols=3,bonds=none]
            \tn[at=(2,2),name=zeroK,skin=box,
                ports={90:virtual:$p$,0:virtual:$q$,270:virtual:$r$,180:virtual:$s$,
                        135:physical:$\sigma$}]{K}
            \tn[at=3 n of zeroK,name=zeroN,skin=box,ports={90:virtual:$a_N$}]{\bra0}
            \tn[at=3 e of zeroK,name=zeroE,skin=box,ports={0:virtual:$a_E$}]{\bra0}
            \tn[at=3 s of zeroK,name=zeroS,skin=box,ports={270:virtual:$a_S$}]{\bra0}
            \tn[at=3 w of zeroK,name=zeroW,skin=box,ports={180:virtual:$a_W$}]{\bra0}
        \end{tenkz}
    \end{tenkzequation}
    Thus this identity connects the original checkerboard tensor to the
    existing color-difference construction, with its coefficient normalization
    unchanged.
\end{theorem}

\begin{proof}\leanok
    A permutation of an orthonormal basis is unitary, and fixing four
    registers at zero takes distinct color basis vectors to distinct boundary
    basis vectors. The coefficient formula of
    Theorem~\ref{thm:peps_kitaev_checkerboard_blocking} gives the matrix
    identity: $E^\dagger$ kills every boundary configuration with a nonzero
    second register and retains the first registers otherwise. The cyclic
    reordering follows from the four differences
    $(p+s,q+p,q+r,r+s)$ of the dual binary tensor.
\end{proof}
